\documentclass{article}
\usepackage{amsmath, amssymb}

\title{1.2.6-60}

\begin{document}
\maketitle

To choose \( k \) elements out of \( n \) without repeatitions.

Let \( X = \{ x_1, x_2, \ldots, x_n \} \)

To choose \( k \) distinct elements from \( X \), we need to get \( x_{a_1}, x_{a_2}, \ldots, x_{a_k} \)

\( 1 \leq a_1 < a_2 < \cdots < a_k \leq n \)

As we know amount of solutions for that inequalities is: \( \binom{n}{k} \)

Now to consider choosing with repeatitions.

\( 1 \leq b_1 \leq b_2 \leq \cdots \leq b_k \leq n \)

\( a_i := b_i + i - 1 \iff b_i = a_i - i + 1 \)

\( b_i \leq b_{i+1} \iff a_i - i + 1 \leq a_{i+1} - (i+1) + 1 \iff a_{i} + 1 \leq a_{i+1} \iff a_i < a_{i+1} \)

\( 1 \leq b_1 \iff 1 \leq a_i \)

\( b_k \leq n \iff a_k - k + 1 \leq n \iff a_k \leq n + k - 1 \)

We converted our inequalities into

\( 1 \leq a_1 < a_2 < \cdots < a_k \leq n - k + 1 \)

So there are \( \binom{n-k+1}{k} \) ways to do so.


\end{document}
